% figures/ch04-q-expansion-geometry.tex
% q 展开的几何意义：τ → q = e^{2πiτ} 把竖带映到穿孔圆盘
\begin{figure}[ht]
\centering
\begin{tikzpicture}[>=Stealth, font=\small, scale=1.1]

% === 左侧：上半平面的竖带 ===
\begin{scope}[xshift=-4.2cm]
  % 阴影填充带
  \fill[blue!8]
    (-1, 0.3) rectangle (1, 3.5);
  % 边界
  \draw[thick, blue!60!black] (-1, 0.3) -- (-1, 3.5);
  \draw[thick, blue!60!black] (1, 0.3) -- (1, 3.5);
  \draw[dashed, gray] (-1.5, 0.3) -- (1.5, 0.3)
    node[right, font=\footnotesize] {$\im(\tau) = Y$};
  % 坐标轴
  \draw[->] (-1.6, 0) -- (1.6, 0) node[right] {$\Re(\tau)$};
  \draw[->] (0, 0) -- (0, 3.8) node[above] {$\im(\tau)$};
  % 刻度
  \node[below] at (-1, 0) {$-\tfrac{1}{2}$};
  \node[below] at (1, 0) {$\tfrac{1}{2}$};
  % 示例点
  \fill[red!70!black] (0.3, 1.5) circle (2pt);
  \node[right, red!70!black, font=\footnotesize] at (0.33, 1.5) {$\tau$};
  % 顶部箭头（∞方向）
  \draw[->, green!50!black, line width=1.2pt] (0, 3.2) -- (0, 3.7);
  \node[green!50!black, right, font=\footnotesize] at (0.05, 3.5) {$\infty$};
  % 标签
  \node[blue!60!black] at (0, 0.9) {$\HH$（竖带）};
\end{scope}

% === 中间箭头 ===
\draw[->, thick, purple!70!black, line width=1.5pt]
  (-1.8, 1.8) -- (-0.3, 1.8)
  node[above, midway, font=\footnotesize, purple!70!black]
  {$q = e^{2\pi i\tau}$};

% === 右侧：穿孔圆盘 ===
\begin{scope}[xshift=1.8cm]
  % 圆盘
  \fill[orange!10] (0,0) circle (2cm);
  \fill[white] (0,0) circle (0.25cm);  % 穿孔
  \draw[thick, orange!70!black] (0,0) circle (2cm);
  \draw[thick, dashed, orange!40!black] (0,0) circle (0.25cm);
  % r=|q|=e^{-2πY}的圆
  \draw[dashed, blue!50!black] (0,0) circle (1.2cm);
  \node[blue!50!black, font=\footnotesize] at (1.05, 0.65)
    {$|q|=e^{-2\pi Y}$};
  % 示例点 q
  \fill[red!70!black] (0.6, 0.9) circle (2pt);
  \node[right, red!70!black, font=\footnotesize] at (0.63, 0.9) {$q(\tau)$};
  % 穿孔标签
  \node[font=\footnotesize, gray] at (0, -0.45) {$q=0$};
  \draw[->, gray, shorten >=3pt] (0.0, -0.42) -- (0.0, -0.23);
  % 外圆标签
  \node[font=\footnotesize, orange!70!black] at (1.6, -1.6)
    {$|q|=1$};
  \node[font=\footnotesize] at (0, -2.5) {穿孔圆盘 $\mathbb{D}^*$};
\end{scope}

% === 对应关系注记 ===
\node[font=\footnotesize, align=left, text=gray!70!black] at (0, -1.5)
  {$\im(\tau)\to+\infty$ $\Leftrightarrow$ $q\to 0$\quad；\quad
   $\tau \equiv \tau+1$ $\Leftrightarrow$ $q$ 不变（周期性）};

\end{tikzpicture}
\caption{$q$-展开的几何意义。映射 $q = e^{2\pi i\tau}$ 把上半平面中的竖带
$\{|\Re(\tau)| \leq \frac{1}{2},\, \im(\tau) > Y\}$（左）
双全纯地映到穿孔圆盘 $\{0 < |q| < e^{-2\pi Y}\}$（右）。
$\tau \to \infty$ 对应 $q \to 0$，模形式在尖点全纯等价于在 $q=0$ 处可去奇点。}
\label{fig:q-expansion-geometry}
\end{figure}
